\begin{textsample}{NEG Dim 6 – sp – Score 1.00 – sp/sp\_inf\_avaliacao.txt}  \label{ex:f6_neg_206}
Processo de avaliação a partir da documentação pedagógica Os Parâmetros Nacionais de Qualidade da Educação Infantil ( BRASIL, 2006 ) explicitam que as experiências vividas em contextos individuais e coletivos constituem -se em importantes informações sobre as crianças, seu desenvolvimento, sua aprendizagem, seus interesses, suas forças e necessidades e precisam ser documentadas, refletidas e compartilhadas com os pais ou responsáveis.
No que se refere ao trabalho dos professores, cabe a eles utilizarem diversos registros, realizados por adultos e crianças, tais como relatórios, fotografias, filmagens, produções infantis, diários, portfólios, murais, dentre outros.
Tais registros servem como instrumento de reflexão sobre as práticas planejadas, na busca de melhores caminhos para acompanhar a aprendizagem e o desenvolvimento da criança.
Assim, a documentação pedagógica deve servir como termômetro para ampliar o olhar e a escuta dos professores com base no contexto da aprendizagem e nas propostas realizadas pelas crianças, historicizando suas vivências e experiências, de forma individual e coletiva, validando o desenvolvimento de suas competências e revelando memórias do seu protagonismo.
O planejamento e a avaliação a partir da documentação pedagógica demandam envolvimento e participação ativa das crianças e dos professores.
As produções infantis, seus pensamentos, interesses, ideias, descobertas, aprendizados, criações, experiências e brincadeiras nos revelam sua maneira de compreender o mundo.
Nesse sentido, os professores precisam registrar as experiências das crianças ( desenhos, produções de textos orais ou escritos, dramatização, momentos da alimentação, dos cuidados de banho e troca etc. ) por meio de filmagens, fotos, portfólios, entre outros, de modo que possam compartilhar os vários saberes com seus pares e com os adultos.
No que se refere à avaliação na Educação Infantil, como já foi dito, esta deve ser realizada por meio de observações e registros, não devendo existir práticas de verificação de aprendizagem tais como as provinhas.
A Lei de Diretrizes e Bases da Educação, na seção II, referente à Educação Infantil, artigo 31, ressalta que : “ [ … ] a avaliação far -se-á mediante o acompanhamento e registro do seu desenvolvimento, sem o objetivo de promoção, mesmo para o acesso ao Ensino Fundamental ”.
No contexto do Currículo Paulista, a documentação pedagógica deve ser vista como um importante instrumento aliado à efetivação da Proposta Pedagógica de cada instituição, ressaltando que aquilo que se documenta e o modo como isso é feito revelam a visão dos sujeitos e as concepções sobre a criança e a escola de educação infantil.

% matched lemmas: 
\end{textsample}
